\scanpage{61}
\uline{Système de Pseudodroites}

\uline{Définition} : c'est une carte (topologique ou combinatoire)
dont la surface sousjacente $X$ est un plan projectif\textsuperscript{*}
% margin note written sideways
\margin{* i.e. (compacte) non orientable et de car. EP égale à 1\ldots}
réel, dont les ordres des sommets sont pairs $\geq4$ (d'où
la division de l'ensemble des arêtes en classes
d'équivalence, appelées les ``pseudodroites'' de la
carte) telles que deux pseudodroites se rencontrent
en un point et un seul.

\uline{Premières conséquences} : le nb \inserted{\deleted{\unsure{$p$}} $n$} des pseudodroites est $\geq2$.
L'immersion de $\mathbb{S}^1$ \unsure{ds|de|dans} $\mathbb{P}_{\mathbb{R}}^2$ correspondant à une
ps. dr. est (des deux classes d'isotopie possibles) celle
qui est non homotope à \deleted{\unsure{un}} $0$, i.e. le long d'une
psdr. $X$ est non orientable, i.e. la situation est la
\unsure{m\textsuperscript{e}|$=$} que pour une \uline{droite} (projective) dans le plan
projectif. On montre que pour des psdr ``en
position générale'' (sommets tous d'ordre 4) et
en nb $n\leq6$, la situation peut être \unsure{réalisé|réalisée} par
un système de vraies droites dans le plan
projectif réel.

Soit $D$ une pseudodroite de $X$, $\widetilde X$ le revêtement des
orientations de $X$ ($\widetilde X\simeq\mathbb{S}^2\simeq\mathbb{P}_{\mathbb{C}}^1$), $\widetilde D$ l'image inverse
de $D$ dans $X$. Alors $\widetilde X\setminus\widetilde D$ a deux comp.
connexes (ou ``hémisphères'') en corr. 1-1 can.
avec l'ens. des deux orientations de $D$.

Notons que \inserted{sur} l'ens. des sommets se trouvant
% archive sheet number: 60
